\scanpage{21}
\mspage{10}
aller comme une brute et de définir
une catégorie plus grosse que la cat.
de Galois-Teichmüller (i.e. avec
de flèches plus gros) par \uline{générateurs}
et relations convenables. Mais alors il
devient problématique, quand le corps de
base est le corps $\mathbb C$ des complexes, d'
interpréter cette catégorie de Galois-Teich-
müller comme une ``complétée profinie''
de la catégorie isotopique envisagée plus
haut.
\margin{20}
